\documentclass[14pt]{extarticle}
\usepackage{fontspec}
\usepackage{polyglossia}
\usepackage{amsmath}
\usepackage{amssymb}
\usepackage{xcolor}
\usepackage{fancyhdr}
\usepackage{graphicx}
\usepackage{listings}
\usepackage[most]{tcolorbox}
% \usepackage{setspace}
\usepackage{pifont}
\usepackage{enumitem}
\usepackage{titlesec}
\usepackage[bottom]{footmisc}
\usepackage{titling}
\usepackage{minted}
\usepackage{etoolbox}
\usepackage{array}
\usepackage{longtable} % page-breakable tables (must load before polyglossia/bidi RTL setup)
\usepackage{extsizes}
\usepackage{float}
\usepackage{booktabs}
\usepackage{hyperref}
\usepackage[onehalfspacing]{setspace}
\newif\ifwhiteblack
\ifdefined\setwhiteblack
  \setwhiteblack
\fi

\ifwhiteblack
  \usepackage[a4paper, landscape, top=1.25cm, bottom=1.25cm, left=1.25cm, right=1.25cm, headheight=1.5cm, headsep=0.5cm]{geometry}
\else
  \usepackage[a4paper, top=2.7cm, bottom=2.7cm, left=2.5cm, right=2.5cm, headheight=1.5cm, headsep=0.5cm]{geometry}
\fi
\usepackage{tikz}
\usepackage{datetime2}

\usetikzlibrary{automata, positioning, arrows.meta, arrows}
\providecommand{\LTcaptype}{table}

\newfontfamily\emoji{Segoe UI Emoji}

\pagestyle{fancy}

\setmainlanguage[numerals=western]{arabic}
\setotherlanguage{english}
\newfontfamily\arabicfont[Script=Arabic]{Calibri}
\newfontfamily\arabicfonttt[Script=Arabic]{Courier New}



\ifwhiteblack
\lstset{
  language=[Sharp]C,
  numbers=left,
  stepnumber=1,
  numbersep=8pt,
  frame=single,
  basicstyle=\ttfamily\small,
  keywordstyle=\bfseries,
  stringstyle=\ttfamily,
  commentstyle=\itshape
}
\else
\lstset{
  language=[Sharp]C,
  numbers=left,
  stepnumber=1,
  numbersep=8pt,
  frame=single,
  basicstyle=\ttfamily\small,
  keywordstyle=\color{blue},
  stringstyle=\color{red},
  commentstyle=\color{green!50!black}
}
\fi

\newif\ifdetailed
\ifdefined\setdetailed
  \setdetailed
\fi

\newif\ifwithsols
\ifdefined\setwithsols
  \setwithsols
\fi

\ifwhiteblack
  \AtBeginDocument{
    \LARGE
    \usemintedstyle{bw}
  }
\fi

% Default minted look: compact code that wraps long lines.
% Not applied in black-and-white landscape mode, which already uses \LARGE.
\ifwhiteblack\else
  \setminted{fontsize=\small, breaklines=true}
\fi

% unified tcolorboxes for articles
\ifwhiteblack
\tcbset{colback=white, colframe=black, fonttitle=\bfseries, boxrule=1pt}
\newtcolorbox{boxDef}[1][]{title={{\emoji📘} تعريف\ifx\\#1\\\else ~#1\fi :}}
\newtcolorbox{boxExercise}[1][]{title={{\emoji🧩} تمرين\ifx\\#1\\\else ~#1\fi :}}
\newtcolorbox{boxExample}[1][]{title={{\emoji📝} مثال\ifx\\#1\\\else ~#1\fi :}}
\newtcolorbox{boxNote}[1][]{title={{\emoji✨} ملاحظة\ifx\\#1\\\else ~#1\fi :}}
\newtcolorbox{boxAttention}[1][]{title={{\emoji🔔} تنبيه\ifx\\#1\\\else ~#1\fi :}}
\newtcolorbox{boxWarning}[1][]{title={{\emoji⚡} ملاحظة هامة\ifx\\#1\\\else ~#1\fi :}}
\newtcolorbox{boxSolution}[1][]{title={{\emoji✅} حل\ifx\\#1\\\else ~#1\fi :}}
\newtcolorbox{boxSymbol}[1][]{title={{\emoji🔣} رمز\ifx\\#1\\\else ~#1\fi :}}
\newtcolorbox{boxHint}[1][]{title={{\emoji💡} تلميح\ifx\\#1\\\else ~#1\fi :}}
\else
\tcbset{colback=white, colframe=black, fonttitle=\bfseries, boxrule=0.8pt}
\newtcolorbox{boxDef}[1][]{colback=blue!5!white,colframe=blue!75!black,
  title={{\emoji📘} تعريف\ifx\\#1\\\else ~#1\fi :}}
\newtcolorbox{boxExercise}[1][]{colback=cyan!5!white,colframe=cyan!70!black,
  title={{\emoji🧩} تمرين\ifx\\#1\\\else ~#1\fi :}}
\newtcolorbox{boxExample}[1][]{colback=yellow!5!white,colframe=orange!90!black,
  title={{\emoji📝} مثال\ifx\\#1\\\else ~#1\fi :}}
\newtcolorbox{boxNote}[1][]{colback=gray!10!white,colframe=black,
  title={{\emoji✨} ملاحظة\ifx\\#1\\\else ~#1\fi :}}
\newtcolorbox{boxAttention}[1][]{colback=magenta!10!white,colframe=magenta!80!black,
  title={{\emoji🔔} تنبيه\ifx\\#1\\\else ~#1\fi :}}
\newtcolorbox{boxWarning}[1][]{colback=red!5!white,colframe=red!75!black,
  title={{\emoji⚡} ملاحظة هامة\ifx\\#1\\\else ~#1\fi :}}
\newtcolorbox{boxSolution}[1][]{colback=green!5!white,colframe=green!60!black,
  title={{\emoji✅} حل\ifx\\#1\\\else ~#1\fi :}}
\newtcolorbox{boxSymbol}[1][]{colback=purple!5!white,colframe=purple!70!black,
  title={{\emoji🔣} رمز\ifx\\#1\\\else ~#1\fi :}}
\newtcolorbox{boxHint}[1][]{colback=teal!5!white,colframe=teal!60!black,
  title={{\emoji💡} تلميح\ifx\\#1\\\else ~#1\fi :}}
\fi


\ifwhiteblack
\tcbset{simplecode/.style={ colback=white, colframe=black, boxrule=1pt, arc=2pt, left=4pt,right=4pt,top=4pt,bottom=4pt}}
\else
\tcbset{simplecode/.style={ colback=gray!5, colframe=black!50, boxrule=0.4pt, arc=2pt, left=4pt,right=4pt,top=4pt,bottom=4pt}}
\fi
% Compact code boxes: tighter line spacing (the preamble uses onehalfspacing) and small vertical gaps.
% Must come AFTER the simplecode definition above.
\tcbset{simplecode/.append style={fontupper=\setstretch{0.95}, top=4pt, bottom=4pt, before skip=6pt, after skip=6pt}}
\newenvironment{boxCode}{\begin{tcolorbox}[simplecode]}{\end{tcolorbox}}

% Code cell used inside tables (keeps code left-to-right inside RTL tables).
% Use inside a fixed-width p{...} column: \begin{pycell}\begin{minted}...\end{minted}\end{pycell}
\newtcolorbox{pycell}{frame empty,colback=white,boxsep=0pt,left=2pt,right=2pt,top=2pt,bottom=2pt,width=\linewidth}

\newcolumntype{C}[1]{>{\centering\arraybackslash}p{#1}}

\makeatletter
\DTMnewdatestyle{dotted}{%
  \renewcommand{\DTMdisplaydate}[4]{\two@digits{##3}.\two@digits{##2}.##1}%
  \renewcommand{\DTMDisplaydate}{\DTMdisplaydate}%
}
\makeatother
\DTMsetdatestyle{dotted}

% redefine spaces after titles
\makeatletter
\renewcommand{\@maketitle}{%
  \begin{center}
    {\huge \bfseries \@title \par}%
    \vskip 0.2em % space between title and author
    {\large \@author \par}%
    % \vskip 0.2em % space between author and date
    % {\normalsize \@date \par}%
    {\small \hfill \begin{english}\DTMtoday\end{english} \par}%
  \end{center}
}
\makeatother

\fancyhf{} % clear default
\fancypagestyle{plain}{
  \fancyhf{}
  \fancyhead[L]{مدرسة التسامح الشاملة}
  % \fancyhead[L]{\includegraphics[height=1cm]{../../../images/logoTasamoh.png}}
  \fancyhead[R]{الأستاذ محمود اغبارية}
  \fancyfoot[C]{\thepage}
}

\fancyhead[L]{مدرسة التسامح الشاملة}
\fancyhead[R]{الأستاذ محمود اغبارية}
\fancyfoot[C]{\thepage}

\setcounter{tocdepth}{3} % only section subsection and subsubsection in TOC

\newcommand{\importpdfpage}[2]{%
    \thispagestyle{empty}
    \newgeometry{left=0cm,right=0cm,top=0cm,bottom=0cm}
    \includegraphics[page=#2,width=0.95\paperwidth,height=0.95\paperheight,keepaspectratio]{#1}
    \restoregeometry
}

\newcommand{\insertFullImg}[1]{%
    \noindent
    \makebox[\textwidth][c]{
        \includegraphics[width=0.9\paperwidth,keepaspectratio]{#1}%
    }%
}%


\ifwithsols
\title{حل مجمّع أسئلة إضافية للاختبار الشهري الأوّل - المتغيرات وجملة الشرط}
\else
\title{مجمّع أسئلة إضافية للاختبار الشهري الأوّل - المتغيرات وجملة الشرط}
\fi

\begin{document}
\maketitle
\thispagestyle{fancy}

\section*{القسم الأوّل: أسئلة فهم وتتبّع كود}

\begin{enumerate}[itemsep=2.2em]

% ==============================================================
% 1. تتبع كود: عدد ثنائي المنازل
% ==============================================================
\item
أمامك مقطع البرنامج التالي الذي يستقبل عددًا صحيحًا موجبًا \textenglish{num} مكوّنًا من منزلتين (بين 10 و 99):

\begin{boxCode}
\begin{english}
\begin{minted}{csharp}
int num = int.Parse(Console.ReadLine());
int a = num / 10;
int b = num % 10;

if (a > b)
{
    if (a % 2 == 0 && b % 2 == 0)
    {
        Console.WriteLine("A");
    }
    else
    {
        Console.WriteLine("B");
    }
}
else
{
    if (a == b || b - a > 4)
    {
        Console.WriteLine("C");
    }
    else
    {
        Console.WriteLine("D");
    }
}
\end{minted}
\end{english}
\end{boxCode}

\begin{enumerate}
    \item \textbf{ما هو مُخرج هذا المُدخل؟} \\
    تتبّع تنفيذ المقطع البرمجي أعلاه لكل قيمة من قيم المُدخل \textenglish{num} في الجدول التالي، واكتب قيمة كلّ من \textenglish{a} و \textenglish{b} والمُخرج الذي يطبعه البرنامج:
    \begin{center}
    \begin{english}
    \renewcommand{\arraystretch}{1.5}
    \begin{tabular}{|C{2.5cm}|C{2.5cm}|C{2.5cm}|C{3.5cm}|}
    \hline
    \textbf{num} & \textbf{a} & \textbf{b} & \textarabic{\textbf{المُخرج}} \\
    \hline
    84 & & & \\
    \hline
    72 & & & \\
    \hline
    35 & & & \\
    \hline
    28 & & & \\
    \hline
    \end{tabular}
    \end{english}
    \end{center}
    \filbreak

    \item \textbf{أعطِ مُدخلاً يطبع كلّ واحد من هذه المخرجات:} \\
    لكل مُخرج في الجدول التالي، اقترح قيمة لعدد ثنائي المنازل \textenglish{num} (\textbf{مختلفة} عن الأعداد الواردة في الفرع أ) تجعل البرنامج يطبع هذا المُخرج، مع تحقيق الشرط الإضافي المذكور:
    \begin{center}
    \renewcommand{\arraystretch}{1.5}
    \begin{tabular}{|C{2.6cm}|p{5.8cm}|C{3.2cm}|}
    \hline
    \textbf{المُخرج المطلوب} & \textbf{شرط إضافي على المُدخل \textenglish{num}} & \textbf{قيمة مقترحة لـ \textenglish{num}} \\
    \hline
    \textenglish{A} & بدون شرط إضافي & \\
    \hline
    \textenglish{B} & بحيث تكون منزلة العشرات \textenglish{a} زوجية & \\
    \hline
    \textenglish{C} & بحيث تكون المنزلتان متساويتين (\textenglish{a == b}) & \\
    \hline
    \textenglish{C} & بحيث تكون المنزلتان مختلفتين (\textenglish{a != b}) & \\
    \hline
    \textenglish{D} & بدون شرط إضافي & \\
    \hline
    \end{tabular}
    \end{center}
    \filbreak
\end{enumerate}

\ifwithsols
\begin{boxSolution}[الفرع أ والفرع ب]
\textbf{حل الفرع أ:}
\begin{center}
\begin{english}
\renewcommand{\arraystretch}{1.4}
\begin{tabular}{|C{2.5cm}|C{2.5cm}|C{2.5cm}|C{3.5cm}|}
\hline
\textbf{num} & \textbf{a} & \textbf{b} & \textarabic{\textbf{المُخرج}} \\
\hline
84 & 8 & 4 & A \\
\hline
72 & 7 & 2 & B \\
\hline
35 & 3 & 5 & D \\
\hline
28 & 2 & 8 & C \\
\hline
\end{tabular}
\end{english}
\end{center}

\textbf{حل الفرع ب (أمثلة لإجابات صحيحة):}
\begin{itemize}[nosep]
    \item لطباعة \textenglish{A}: يجب أن يكون \textenglish{a > b} وكلاهما زوجي، مثل: \textenglish{62} أو \textenglish{64} أو \textenglish{82}.
    \item لطباعة \textenglish{B} مع كون \textenglish{a} زوجيًا: يجب أن يكون \textenglish{a > b} و \textenglish{b} فرديًا، مثل: \textenglish{63} أو \textenglish{85} أو \textenglish{41}.
    \item لطباعة \textenglish{C} مع تساوي المنزلتين: مثل \textenglish{44} أو \textenglish{77}.
    \item لطباعة \textenglish{C} مع اختلاف المنزلتين: يجب أن يكون \textenglish{b - a > 4}، مثل \textenglish{17} أو \textenglish{39}.
    \item لطباعة \textenglish{D}: يجب أن يكون \textenglish{a < b} والفرق بينهما $4$ أو أقل، مثل \textenglish{25} أو \textenglish{47}.
\end{itemize}
\end{boxSolution}
\fi
\filbreak

% ==============================================================
% 2. تتبع كود: 3 أعداد وعامل النفي
% ==============================================================
\item
أمامك مقطع البرنامج التالي الذي يستقبل ثلاثة أعداد صحيحة \textenglish{a} و \textenglish{b} و \textenglish{c}:

\begin{boxCode}
\begin{english}
\begin{minted}{csharp}
int a = int.Parse(Console.ReadLine());
int b = int.Parse(Console.ReadLine());
int c = int.Parse(Console.ReadLine());

if (a + b > c)
{
    if (a == b && b != c)
    {
        Console.WriteLine("A");
    }
    else if (c % 2 == 0)
    {
        Console.WriteLine("B");
    }
    else
    {
        Console.WriteLine("C");
    }
}
else
{
    if (!(a > 0 && b > 0))
    {
        Console.WriteLine("D");
    }
    else
    {
        Console.WriteLine("E");
    }
}
\end{minted}
\end{english}
\end{boxCode}

\begin{enumerate}
    \item \textbf{ما هو مُخرج هذا المُدخل؟} \\
    تتبّع تنفيذ المقطع البرمجي أعلاه لكل ثلاثية من المدخلات في الجدول التالي، واكتب المُخرج الذي يطبعه البرنامج:
    \begin{center}
    \begin{english}
    \renewcommand{\arraystretch}{1.5}
    \begin{tabular}{|C{2.2cm}|C{2.2cm}|C{2.2cm}|C{3.8cm}|}
    \hline
    \textbf{a} & \textbf{b} & \textbf{c} & \textarabic{\textbf{المُخرج}} \\
    \hline
    5 & 5 & 8 & \\
    \hline
    4 & 4 & 4 & \\
    \hline
    6 & 3 & 7 & \\
    \hline
    2 & 3 & 9 & \\
    \hline
    \end{tabular}
    \end{english}
    \end{center}
    \filbreak

    \item \textbf{أعطِ مُدخلاً يطبع كلّ واحد من هذه المخرجات:} \\
    لكل مُخرج في الجدول التالي، اقترح قيَمًا للأعداد الصحيحة \textenglish{a} و \textenglish{b} و \textenglish{c} (\textbf{مختلفة} عن ثلاثيات الفرع أ) تجعل البرنامج يطبع هذا المُخرج، مع تحقيق الشرط الإضافي المذكور:
    \begin{center}
    \renewcommand{\arraystretch}{1.5}
    \begin{tabular}{|C{2.2cm}|p{4.8cm}|C{1.5cm}|C{1.5cm}|C{1.5cm}|}
    \hline
    \textbf{المُخرج المطلوب} & \textbf{شرط إضافي على المُدخلات} & \textbf{\textenglish{a}} & \textbf{\textenglish{b}} & \textbf{\textenglish{c}} \\
    \hline
    \textenglish{A} & بدون شرط إضافي & & & \\
    \hline
    \textenglish{B} & بحيث يكون \textenglish{a != b} & & & \\
    \hline
    \textenglish{C} & بحيث يكون \textenglish{a == b} & & & \\
    \hline
    \textenglish{D} & بحيث يكون \textenglish{a > 0} & & & \\
    \hline
    \textenglish{E} & بدون شرط إضافي & & & \\
    \hline
    \end{tabular}
    \end{center}
    \filbreak
\end{enumerate}

\ifwithsols
\begin{boxSolution}[الفرع أ والفرع ب]
\textbf{حل الفرع أ:}
\begin{center}
\begin{english}
\renewcommand{\arraystretch}{1.4}
\begin{tabular}{|C{2.2cm}|C{2.2cm}|C{2.2cm}|C{3.8cm}|}
\hline
\textbf{a} & \textbf{b} & \textbf{c} & \textarabic{\textbf{المُخرج}} \\
\hline
5 & 5 & 8 & A \\
\hline
4 & 4 & 4 & B \\
\hline
6 & 3 & 7 & C \\
\hline
2 & 3 & 9 & E \\
\hline
\end{tabular}
\end{english}
\end{center}

\textbf{حل الفرع ب (أمثلة لإجابات صحيحة):}
\begin{itemize}[nosep]
    \item لطباعة \textenglish{A}: يجب أن يكون \textenglish{a + b > c} و \textenglish{a == b} و \textenglish{b != c}، مثل: \textenglish{a = 3, b = 3, c = 4}.
    \item لطباعة \textenglish{B} مع \textenglish{a != b}: يجب أن يكون \textenglish{a + b > c} و \textenglish{c} زوجيًا، مثل: \textenglish{a = 5, b = 4, c = 6}.
    \item لطباعة \textenglish{C} مع \textenglish{a == b}: يجب أن تكون الأعداد الثلاثة متساوية وفردية (\textenglish{a == b == c})، مثل: \textenglish{a = 3, b = 3, c = 3}.
    \item لطباعة \textenglish{D} مع \textenglish{a > 0}: يجب أن يكون \textenglish{b <= 0} و \textenglish{a + b <= c}، مثل: \textenglish{a = 4, b = -2, c = 5}.
    \item لطباعة \textenglish{E}: يجب أن يكون \textenglish{a > 0} و \textenglish{b > 0} و \textenglish{a + b <= c}، مثل: \textenglish{a = 1, b = 2, c = 5}.
\end{itemize}
\end{boxSolution}
\fi
\filbreak

\end{enumerate}

\clearpage
\section*{القسم الثاني: أسئلة كتابة برنامج}

\begin{enumerate}[resume, itemsep=2.2em]

% ==============================================================
% 3. حذف المنزلة الوسطى وفحص القسمة
% ==============================================================
\item
اكتب برنامجًا يستقبل من المستخدم عددًا صحيحًا موجبًا مكوّنًا من \textbf{3 منازل} (بين 100 و 999). \\
يقوم البرنامج بتكوين \textbf{عدد جديد مكوّن من منزلتين} ناتج عن حذف منزلة العشرات الوسطى من العدد الأصلي (فمثلاً من العدد \textenglish{632} يتكوّن العدد الجديد \textenglish{62})، ويطبع هذا العدد الجديد. \\
بعد ذلك، يفحص البرنامج العلاقة بين العدد الجديد ومنزلة العشرات التي حُذفت، ويطبع إحدى الرسائل التالية:
\begin{itemize}[nosep]
    \item إذا كانت منزلة العشرات المحذوفة تساوي صفرًا، لا يمكن القسمة عليها، فيطبع البرنامج: \textenglish{"Middle digit is zero"}.
    \item وإلا، إذا كان العدد الجديد يقبل القسمة على منزلة العشرات المحذوفة بدون باقٍ، يطبع البرنامج: \textenglish{"Divisible"}.
    \item وإلا، يطبع البرنامج: \textenglish{"Not divisible"}.
\end{itemize}

\ifdetailed
\begin{boxExample}[1]
\begin{english}
\begin{minted}{text}
Enter a 3-digit number:
633
63
Divisible
\end{minted}
\end{english}
\end{boxExample}
لأنّ بعد حذف منزلة العشرات ($3$) من العدد $633$ يتكوّن العدد الجديد $63$، وبما أنّ $63$ يقبل القسمة على $3$ بدون باقٍ، يطبع البرنامج \textenglish{Divisible}.
\begin{boxExample}[2]
\begin{english}
\begin{minted}{text}
Enter a 3-digit number:
408
48
Middle digit is zero
\end{minted}
\end{english}
\end{boxExample}
\fi

\ifwithsols
\begin{boxSolution}
\begin{english}
\begin{minted}{csharp}
Console.WriteLine("Enter a 3-digit number:");
int num = int.Parse(Console.ReadLine());

int hundreds = num / 100;
int tens = (num / 10) % 10;
int units = num % 10;

int newNum = hundreds * 10 + units;
Console.WriteLine(newNum);

if (tens == 0)
{
    Console.WriteLine("Middle digit is zero");
}
else if (newNum % tens == 0)
{
    Console.WriteLine("Divisible");
}
else
{
    Console.WriteLine("Not divisible");
}
\end{minted}
\end{english}
\end{boxSolution}
\fi
\filbreak

% ==============================================================
% 4. حافلات الرحلة المدرسية
% ==============================================================
\item
تنظّم المدرسة رحلة تعليمية، وتريد حجز حافلات لنقل الطلاب. تتسع كل \textbf{حافلة كبيرة} لـ \textenglish{45} طالبًا بالضبط. \\
اكتب برنامجًا يستقبل من المستخدم العدد الكلي للطلاب المسجلين للرحلة (عدد صحيح موجب). \\
يطبع البرنامج أولاً عدد \textbf{الحافلات الكبيرة الممتلئة بالكامل}، ثم عدد \textbf{الطلاب المتبقين} الذين لا يكفون لملء حافلة كبيرة كاملة. \\
بعد ذلك، يطبع البرنامج توصية بخصوص الطلاب المتبقين حسب القواعد التالية:
\begin{itemize}[nosep]
    \item إذا لم يتبقَّ أي طالب (الباقي يساوي $0$)، يطبع: \textenglish{"No extra transport needed"}.
    \item إذا كان عدد الطلاب المتبقين بين $1$ و $15$ طالبًا (شاملاً $1$ و $15$)، يطبع: \textenglish{"Need 1 minibus"}.
    \item إذا كان عدد الطلاب المتبقين أكثر من $15$ طالبًا، يطبع: \textenglish{"Need 1 extra large bus"}.
\end{itemize}

\ifdetailed
\begin{boxExample}[1]
\begin{english}
\begin{minted}{text}
Enter total students:
100
Full large buses: 2
Remaining students: 10
Need 1 minibus
\end{minted}
\end{english}
\end{boxExample}
\begin{boxExample}[2]
\begin{english}
\begin{minted}{text}
Enter total students:
115
Full large buses: 2
Remaining students: 25
Need 1 extra large bus
\end{minted}
\end{english}
\end{boxExample}
\fi

\ifwithsols
\begin{boxSolution}
\begin{english}
\begin{minted}{csharp}
Console.WriteLine("Enter total students:");
int students = int.Parse(Console.ReadLine());

int fullBuses = students / 45;
int remaining = students % 45;

Console.WriteLine("Full large buses: " + fullBuses);
Console.WriteLine("Remaining students: " + remaining);

if (remaining == 0)
{
    Console.WriteLine("No extra transport needed");
}
else if (remaining <= 15)
{
    Console.WriteLine("Need 1 minibus");
}
else
{
    Console.WriteLine("Need 1 extra large bus");
}
\end{minted}
\end{english}
\end{boxSolution}
\fi
\filbreak

% ==============================================================
% 5. الطاولة الدائرية داخل الغرفة المربعة
% ==============================================================
\item
يريد مهندس ديكور فحص هل تتسع طاولة دائرية الشكل داخل غرفة مربّعة الشكل. \\
اكتب برنامجًا يستقبل من المستخدم مساحة الغرفة المربّعة بالمتر المربّع (عدد عشري موجب)، ثم يستقبل طول نصف قطر الطاولة الدائرية بالمتر (عدد عشري موجب). \\
يحسب البرنامج ويطبع \textbf{طول ضلع الغرفة المربّعة} ($\text{side} = \sqrt{\text{roomArea}}$)، وكذلك \textbf{قُطر الطاولة الدائرية} ($\text{diameter} = 2 \times \text{tableRadius}$)، ثم يطبع إحدى الرسائل التالية:
\begin{itemize}[nosep]
    \item إذا كان طول ضلع الغرفة أكبر من قُطر الطاولة بـ \textbf{مترين على الأقل} (أي الفرق بينهما $2$ أو أكثر)، يطبع: \textenglish{"Fits comfortably"}.
    \item إذا كان طول ضلع الغرفة أكبر من أو يساوي قُطر الطاولة، ولكن الفرق بينهما أقل من مترين، يطبع: \textenglish{"Fits tightly"}.
    \item إذا كان قُطر الطاولة أكبر من طول ضلع الغرفة، يطبع: \textenglish{"Does not fit"}.
\end{itemize}

\ifdetailed
\begin{boxExample}[1]
\begin{english}
\begin{minted}{text}
Enter square room area:
36
Enter table radius:
1.5
Room side: 6
Table diameter: 3
Fits comfortably
\end{minted}
\end{english}
\end{boxExample}
لأنّ ضلع الغرفة هو $\sqrt{36} = 6$ أمتار، وقُطر الطاولة هو $2 \times 1.5 = 3$ أمتار، والفرق بينهما $6 - 3 = 3 \ge 2$.
\fi

\ifwithsols
\begin{boxSolution}
\begin{english}
\begin{minted}{csharp}
Console.WriteLine("Enter square room area:");
double roomArea = double.Parse(Console.ReadLine());
Console.WriteLine("Enter table radius:");
double tableRadius = double.Parse(Console.ReadLine());

double side = Math.Sqrt(roomArea);
double diameter = 2 * tableRadius;

Console.WriteLine("Room side: " + side);
Console.WriteLine("Table diameter: " + diameter);

if (side - diameter >= 2)
{
    Console.WriteLine("Fits comfortably");
}
else if (side >= diameter)
{
    Console.WriteLine("Fits tightly");
}
else
{
    Console.WriteLine("Does not fit");
}
\end{minted}
\end{english}
\end{boxSolution}
\fi
\filbreak

% ==============================================================
% 6. حملة تخفيضات متجر الكتب
% ==============================================================
\item
أعلن متجر للكتب عن حملة تخفيضات خاصة عند شراء \textbf{3 كتب}. \\
اكتب برنامجًا يستقبل من المستخدم أسعار الكتب الثلاثة (أعداد عشرية موجبة)، ثم يستقبل قيمة منطقية (\textenglish{bool}) تمثّل هل الزبون مشترك في نادي القرّاء (\textenglish{true} للمشترك و \textenglish{false} لغير المشترك). \\
يحسب البرنامج المبلغ النهائي للدفع حسب القواعد المتسلسلة التالية:
\begin{enumerate}
    \item التكلفة الأولية هي مجموع أسعار الكتب الثلاثة.
    \item \textbf{العرض الأوّل (تقارب الأسعار):} إذا كان الفرق بين سعر \textbf{أغلى كتاب} وسعر \textbf{أرخص كتاب} من بين الكتب الثلاثة \textbf{لا يتجاوز $20$ شيكلاً} (أي أصغر من أو يساوي $20$)، يحصل الزبون على تخفيض بقيمة \textbf{نصف سعر الكتاب الأرخص} (يُطرح نصف سعر الكتاب الأرخص من التكلفة).
    \item \textbf{العرض الثاني (نادي القرّاء والمشتريات الكبيرة):} بعد فحص العرض الأول، إذا كان الزبون مشتركًا في نادي القرّاء \textbf{أو} كانت التكلفة الحالية أكثر من $250$ شيكلاً، يُطرح من التكلفة مبلغ إضافي ثابت قدره $15$ شيكلاً.
    \item في النهاية، يطبع البرنامج المبلغ النهائي للدفع.
\end{enumerate}

\ifdetailed
\begin{boxExample}
\begin{english}
\begin{minted}{text}
Enter 3 book prices:
80
90
70
Is club member (true/false):
true
190
\end{minted}
\end{english}
\end{boxExample}
لأنّ مجموع أسعار الكتب هو $80 + 90 + 70 = 240$. أغلى كتاب هو $90$ وأرخص كتاب هو $70$، والفرق بينهما $90 - 70 = 20 \le 20$، لذلك يحصل على خصم نصف سعر الأرخص ($35$ شيكلاً) فتصبح التكلفة $240 - 35 = 205$. وبما أنه مشترك في النادي (\textenglish{true})، يحصل على خصم $15$ شيكلاً إضافية فيصبح المبلغ النهائي $205 - 15 = 190$.
\fi

\ifwithsols
\begin{boxSolution}
\begin{english}
\begin{minted}{csharp}
Console.WriteLine("Enter 3 book prices:");
double p1 = double.Parse(Console.ReadLine());
double p2 = double.Parse(Console.ReadLine());
double p3 = double.Parse(Console.ReadLine());

Console.WriteLine("Is club member (true/false):");
bool isMember = bool.Parse(Console.ReadLine());

double total = p1 + p2 + p3;

double maxPrice = Math.Max(p1, Math.Max(p2, p3));
double minPrice = Math.Min(p1, Math.Min(p2, p3));

if (maxPrice - minPrice <= 20)
{
    total = total - (minPrice / 2.0);
}

if (isMember || total > 250)
{
    total = total - 15;
}

Console.WriteLine(total);
\end{minted}
\end{english}
\end{boxSolution}
\fi
\filbreak

% ==============================================================
% 7. فحص قوة الرقم السري PIN من 4 منازل
% ==============================================================
\item
يفحص نظام بنكي مدى قوّة الرقم السرّي (\textenglish{PIN}) للبطاقة الائتمانية، وهو عدد صحيح موجب مكوّن من \textbf{4 منازل} (بين 1000 و 9999، مثلاً \textenglish{3839}). \\
اكتب برنامجًا يستقبل من المستخدم عددًا صحيحًا موجبًا من 4 منازل (\textenglish{pin}). \\
يحسب البرنامج ويطبع \textbf{مجموع الخانتين الخارجيتين} (أي مجموع خانة الآلاف وخانة الآحاد)، وكذلك \textbf{مجموع الخانتين الداخليتين} (أي مجموع خانة المئات وخانة العشرات). \\
بعد ذلك، يطبع البرنامج \textenglish{"Strong PIN"} إذا تحقّقت \textbf{جميع} الشروط الثلاثة التالية معًا، وإلا يطبع \textenglish{"Weak PIN"}:
\begin{itemize}[nosep]
    \item خانة الآلاف لا تساوي خانة الآحاد.
    \item الفرق المطلق (الموجب) بين مجموع الخانتين الخارجيتين ومجموع الخانتين الداخليتين يساوي $1$ بالضبط.
    \item أكبر خانة من بين الخانات الأربع أكبر من أو تساوي $7$.
\end{itemize}

\ifdetailed
\begin{boxExample}[1]
\begin{english}
\begin{minted}{text}
Enter 4-digit PIN:
3849
Outer digits sum: 12
Inner digits sum: 12
Weak PIN
\end{minted}
\end{english}
\end{boxExample}
لأنّ في العدد $3849$: مجموع الخارجيتين هو $3 + 9 = 12$ ومجموع الداخليتين هو $8 + 4 = 12$، والفرق المطلق بينهما $0 \neq 1$.
\begin{boxExample}[2]
\begin{english}
\begin{minted}{text}
Enter 4-digit PIN:
3839
Outer digits sum: 12
Inner digits sum: 11
Strong PIN
\end{minted}
\end{english}
\end{boxExample}
لأنّ في العدد $3839$: مجموع الخارجيتين $3 + 9 = 12$، ومجموع الداخليتين $8 + 3 = 11$، والفرق المطلق $|12 - 11| = 1$، والآلاف لا تساوي الآحاد ($3 \neq 9$)، وأكبر خانة هي $9 \ge 7$، فيطبع \textenglish{Strong PIN}.
\fi

\ifwithsols
\begin{boxSolution}
\begin{english}
\begin{minted}{csharp}
Console.WriteLine("Enter 4-digit PIN:");
int pin = int.Parse(Console.ReadLine());

int units = pin % 10;
int tens = (pin / 10) % 10;
int hundreds = (pin / 100) % 10;
int thousands = pin / 1000;

int outerSum = thousands + units;
int innerSum = hundreds + tens;

Console.WriteLine("Outer digits sum: " + outerSum);
Console.WriteLine("Inner digits sum: " + innerSum);

int max1 = Math.Max(thousands, hundreds);
int max2 = Math.Max(tens, units);
int maxDigit = Math.Max(max1, max2);

if (thousands != units && Math.Abs(outerSum - innerSum) == 1 && maxDigit >= 7)
{
    Console.WriteLine("Strong PIN");
}
else
{
    Console.WriteLine("Weak PIN");
}
\end{minted}
\end{english}
\end{boxSolution}
\fi
\filbreak

% ==============================================================
% 8. مراقبة مواعيد القطارات HHMM
% ==============================================================
\item
تراقب شركة القطارات دقة مواعيد وصول قطاراتها. يُمثَّل الوقت كعدد صحيح مكوّن من 4 منازل بنظام 24 ساعة (\textenglish{HHMM}، حيث تمثل المنزلتان اليسريان الساعات والمنزلتان اليمنيان الدقائق؛ مثلاً العدد \textenglish{1420} يعني الساعة \textenglish{14:20}، والساعة الواحدة تساوي $60$ دقيقة). \\
اكتب برنامجًا يستقبل من المستخدم وقت الوصول المجدول (\textenglish{scheduled}) ثم وقت الوصول الفعلي (\textenglish{actual}) لقطار معيّن في نفس اليوم، ويتصرّف حسب القواعد التالية:
\begin{itemize}
    \item إذا وصل القطار في الموعد المجدول أو قبله:
    \begin{itemize}[nosep]
        \item إذا وصل مبكرًا بأكثر من $10$ دقائق، يطبع البرنامج: \textenglish{"Too Early"}.
        \item وإلا، يطبع البرنامج: \textenglish{"On Time"}.
    \end{itemize}
    \item إذا وصل القطار متأخرًا عن الموعد المجدول:
    \begin{itemize}[nosep]
        \item يطبع البرنامج عدد دقائق التأخير (\textenglish{delay}).
        \item يحسب البرنامج نسبة التعويض للمسافرين حسب المعادلة:
        $$ \text{compensation} = \sqrt{\text{delay}} \times 4.5 $$
        ثم يقرّب النتيجة إلى أقرب عدد صحيح ويطبعها.
    \end{itemize}
\end{itemize}

\ifdetailed
\begin{boxExample}[1]
\begin{english}
\begin{minted}{text}
Enter scheduled arrival (HHMM):
1350
Enter actual arrival (HHMM):
1406
Delay in minutes: 16
Compensation: 18
\end{minted}
\end{english}
\end{boxExample}
لأنّ الوقت المجدول هو $13 \times 60 + 50 = 830$ دقيقة، والوقت الفعلي هو $14 \times 60 + 6 = 846$ دقيقة. التأخير هو $846 - 830 = 16$ دقيقة، والتعويض هو $\sqrt{16} \times 4.5 = 4 \times 4.5 = 18$.
\fi

\ifwithsols
\begin{boxSolution}
\begin{english}
\begin{minted}{csharp}
Console.WriteLine("Enter scheduled arrival (HHMM):");
int scheduled = int.Parse(Console.ReadLine());
Console.WriteLine("Enter actual arrival (HHMM):");
int actual = int.Parse(Console.ReadLine());

int schedMins = (scheduled / 100) * 60 + (scheduled % 100);
int actualMins = (actual / 100) * 60 + (actual % 100);

if (actualMins <= schedMins)
{
    if (schedMins - actualMins > 10)
    {
        Console.WriteLine("Too Early");
    }
    else
    {
        Console.WriteLine("On Time");
    }
}
else
{
    int delay = actualMins - schedMins;
    Console.WriteLine("Delay in minutes: " + delay);
    double comp = Math.Round(Math.Sqrt(delay) * 4.5);
    Console.WriteLine("Compensation: " + comp);
}
\end{minted}
\end{english}
\end{boxSolution}
\fi
\filbreak

% ==============================================================
% 9. تكلفة شحن الطرود البريدية
% ==============================================================
\item
تحسب شركة بريد تكلفة شحن الطرود بناءً على منطقة الشحن، وزن الطرد، وما إذا كان الطرد قابلاً للكسر. \\
اكتب برنامجًا يستقبل من المستخدم:
\begin{itemize}[nosep]
    \item حرفًا واحدًا (\textenglish{char}) يمثّل منطقة الشحن: \textenglish{'L'} للشحن المحلي، أو \textenglish{'I'} للشحن الدولي.
    \item عددًا عشريًا موجبًا (\textenglish{double}) يمثّل وزن الطرد بالكيلوغرام.
    \item قيمة منطقية (\textenglish{bool}) تحدّد هل الطرد قابل للكسر (\textenglish{true} إذا كان قابلاً للكسر، و \textenglish{false} وإلا).
\end{itemize}

يحسب البرنامج ويطبع التكلفة النهائية للشحن حسب القواعد التالية:
\begin{itemize}
    \item \textbf{التكلفة الأساسية حسب المنطقة والوزن:}
    \begin{itemize}[nosep]
        \item في حالة الشحن المحلي (\textenglish{'L'}): إذا كان الوزن حتى $5$ كغم (شاملاً $5$) تكون التكلفة الأساسية $25$ شيكلاً؛ وإذا زاد الوزن عن $5$ كغم تكون التكلفة الأساسية $25$ شيكلاً لأول $5$ كغم بالإضافة إلى $4$ شواكل عن كل كيلوغرام إضافي فوق الـ $5$.
        \item في حالة الشحن الدولي (\textenglish{'I'}): إذا كان الوزن حتى $5$ كغم (شاملاً $5$) تكون التكلفة الأساسية $80$ شيكلاً؛ وإذا زاد الوزن عن $5$ كغم تكون التكلفة الأساسية $80$ شيكلاً لأول $5$ كغم بالإضافة إلى $15$ شيكلاً عن كل كيلوغرام إضافي فوق الـ $5$.
    \end{itemize}
    \item \textbf{رسوم التأمين:} إذا كان الطرد قابلاً للكسر، تُضاف إلى التكلفة الأساسية رسوم تأمين بنسبة $10\%$ من التكلفة الأساسية.
\end{itemize}

\ifdetailed
\begin{boxExample}
\begin{english}
\begin{minted}{text}
Enter zone (L or I):
I
Enter weight in kg:
8
Is fragile (true/false):
true
137.5
\end{minted}
\end{english}
\end{boxExample}
لأنّ الشحن دولي (\textenglish{'I'}) والوزن $8 > 5$، فالتكلفة الأساسية هي $80 + (8 - 5) \times 15 = 80 + 45 = 125$ شيكلاً. وبما أن الطرد قابل للكسر (\textenglish{true})، تُضاف $10\%$ ($12.5$ شيكل) فتصبح التكلفة النهائية $125 + 12.5 = 137.5$.
\fi

\ifwithsols
\begin{boxSolution}
\begin{english}
\begin{minted}{csharp}
Console.WriteLine("Enter zone (L or I):");
char zone = char.Parse(Console.ReadLine());
Console.WriteLine("Enter weight in kg:");
double weight = double.Parse(Console.ReadLine());
Console.WriteLine("Is fragile (true/false):");
bool isFragile = bool.Parse(Console.ReadLine());

double cost;

if (zone == 'L')
{
    if (weight <= 5)
    {
        cost = 25;
    }
    else
    {
        cost = 25 + (weight - 5) * 4;
    }
}
else
{
    if (weight <= 5)
    {
        cost = 80;
    }
    else
    {
        cost = 80 + (weight - 5) * 15;
    }
}

if (isFragile)
{
    cost = cost + cost * 0.10;
}

Console.WriteLine(cost);
\end{minted}
\end{english}
\end{boxSolution}
\fi
\filbreak

\end{enumerate}

\vspace{1cm}
\begin{flushleft}
أرجو لكم وقتًا ممتعًا.

الأستاذ محمود اغبارية.
\end{flushleft}

\end{document}
